\documentclass[10pt,a4paper]{article}
\input{../preamble.tex}
%% AUTO-GENERATED by tools/gen_codebase_table.py — do not edit by hand.
\newcommand{\EMFileCount}{189}
\newcommand{\EMLineCount}{96{,}450}
\newcommand{\EMFileCountAll}{208}
\newcommand{\EMLineCountAll}{100{,}791}

%% GENERATED by tools/gen_dead_ends.py
\newcommand{\DEnumbers}{179}
\newcommand{\DEentries}{169}
\newcommand{\DEwitnessed}{37}
\newcommand{\DErevivable}{15}
\newcommand{\DEunassigned}{10}

\providecommand{\CME}{\mathrm{CME}}
\providecommand{\CCSB}{\mathrm{CCSB}}
\newcommand{\lpf}{\operatorname{lpf}}
\newcommand{\gpf}{\operatorname{gpf}}
\newcommand{\RD}{\mathrm{RD}}
\newcommand{\SHH}{\mathrm{SHH}}
\newcommand{\Cinf}{(\mathrm{C}\infty)}

\title{The Euclid--Mullin sequence as an orbit:\\
reformulations, a composite floor, and population laws\\
{\large (with a Lean~4 formalization)}}
\author{Marcello Paris}
\date{August 2026}

\begin{document}
\maketitle

\begin{abstract}
The Euclid--Mullin sequence is $a(0)=2$, $a(n{+}1)=\lpf(a(0)\cdots a(n)+1)$,
$\lpf$ the least prime factor; Mullin's conjecture (MC, 1963) says every
prime appears.  The running product $P_n=a(0)\cdots a(n)$ evolves by the \emph{Euclid step}
$T(n)=n\cdot\lpf(n{+}1)$ on squarefree integers, so the sequence is the
orbit of~$2$ under~$T$ and MC asks whether that orbit meets every prime.  We record what is provable about one orbit of~$T$, all of it
machine-checked in Lean~4 over Mathlib.  (i)~\emph{Reformulations.}
Modulo a prime~$q$ the orbit is a multiplicative walk; $q$ appears iff the
walk hits~$-1$ at the right stage, and MC is equivalent to each of several
walk statements---cofinal hitting, character-sum cancellation, and
\emph{conditional multiplier equidistribution} (CME).  These are
equivalences, not reductions: each is quantified over missing primes, so
MC empties the quantifier; their content is the direction turning a
hitting problem into an equidistribution problem.  (ii)~\emph{The
composite floor.}  $C=\lim\log(a(0)\cdots a(N))/2^{N}$ exists, is a
semiconjugacy $C(Tm)=2C(m)$, and vanishes iff infinitely many Euclid
numbers are composite~$\Cinf$; MC implies $\Cinf$, and $\Cinf$ says
exactly that the sequence is not eventually a Sylvester-type tower of
primes---an open, Fermat-shaped statement.  A second, unscaled invariant
$\rho$ (relative size of the multiplier) satisfies $\rho\in[0,\tfrac12]\cup\{1\}$
and inserts the rung $\mathrm{MC}\Rightarrow\mathrm{RD}\Rightarrow\rho=0\Rightarrow\Cinf$.
(iii)~\emph{Population laws.}  The density of $\{m:\lpf m=p\}$ is
$w_p=p^{-1}\prod_{r<p}(1-1/r)$, the $w_p$ telescope, and the class of
$\lpf\bmod q$ among $q$-rough integers has the \emph{convergent} density
$\sum_{p\equiv a}w_p$; so fixed-modulus equidistribution of the least
prime factor is equivalent to a family of series identities untouched by
Dirichlet's theorem, which fails by head domination.  On
the progression $m\equiv1\pmod P$ containing the Euclid numbers, the least
prime outside the bag is selected with relative density exactly
$1/q$---Shanks' heuristic made precise, and the $n=0$ case of an exact
type-conditional selection law.  Iterating that law
against a deterministic per-path compensator bound yields, for each
\emph{fixed} prime~$q$ and each $\varepsilon>0$, a finite horizon~$n$ at
which the seeds coprime to~$q$ missing~$q$ in their first~$n$ multipliers
have upper natural density at most~$\varepsilon$: unconditional, with no
character sums, and silent about the orbit of~$2$.
Simultaneity in~$q$ is an additivity question, not a rate question: open
for natural density past finite sets of primes, while on the profinite
ensemble $\prod_{r\ \text{prime}}\ZZ/r\ZZ$ with the product uniform measure the
same counting chain gives $\mu$-almost every point capturing every prime
whose coordinate it does not annihilate --- a different ensemble, in which
$\NN$ is null.
For each prime~$q$, almost every
squarefree seed also reaches~$q$ somewhere in its factor tree, unconditionally.
Karamata's Tauberian theorem and Mertens' theorem in progressions are
formalized along the way.  (iv)~\emph{Min versus max.}  Cox--van der
Poorten's omission of~$5$ for the largest-prime-factor variant is
machine-checked; under $\lpf$ capture is a residue-class condition, so no
propagating congruence invariant blocks a prime.
\end{abstract}

\section{Introduction}
\label{sec:intro}

Euclid's proposition~IX.20 of the \emph{Elements} shows that for any finite
set of distinct primes, \emph{every} prime factor of their product plus one
lies outside the set: to grow your set of primes you may pick any factor of
their product plus one
(\href{https://open.substack.com/pub/marcellop71/p/a-subtle-misreading-of-euclid}{a
Substack post I wrote on this}).  Mullin~\cite{Mullin1963} made the choice
canonical: always take the \emph{smallest} prime factor.
\begin{equation}\label{eq:em}
  a(0)=2, \qquad a(n{+}1)=\lpf\bigl(a(0)a(1)\cdots a(n)+1\bigr).
\end{equation}
The first terms are $2,3,7,43,13,53,5,6221671,38709183810571,139,2801,11,
17,5471,\dots$; only 51 terms are known, and small primes such as $41$ and
$47$ have not yet been seen.  Every prime appears at most once, since a
prime dividing $a(0)\cdots a(n)$ cannot divide $a(0)\cdots a(n)+1$.

\begin{conjecture}[Mullin 1963]\label{conj:mc}
Every prime appears in~\eqref{eq:em}.
\end{conjecture}

Write $P_n=a(0)\cdots a(n)$ for the running product.  It records the
whole history---the set of primes collected so far, of which it is the
product---and one step of~\eqref{eq:em} multiplies it by the least prime
factor of $P_n+1$.  So $P_n$ evolves by a map that depends only on its
current value, the \emph{Euclid step}
\[
  T(n)=n\cdot\lpf(n{+}1),\qquad P_{n+1}=T(P_n),\qquad P_n=T^{n}(2),
\]
defined on the squarefree integers: if $n$ is squarefree then $\lpf(n{+}1)$
is a prime not dividing~$n$, so $T(n)$ is squarefree with one more prime
factor and $T(n)\ge2n$.  A squarefree integer \emph{is} a finite set of
primes, and $T$ is Euclid's construction with Mullin's choice, iterated;
$a(n{+}1)=T(P_n)/P_n$ is the prime appended at step~$n$.  Mullin's
conjecture (MC) asks whether one orbit of~$T$---the orbit of~$2$---meets
every prime.  Shanks~\cite{Shanks1991} observed
that modulo a prime~$q$ the orbit is a multiplicative walk on
$(\ZZ/q\ZZ)^\times$ and argued heuristically that a pseudo-random walk
must hit $-1$, after which $q$ is a candidate for selection.  For the
variant that takes the \emph{largest} prime factor, Cox and van der
Poorten~\cite{CoxVdP1968} proved that $5,11,13,\dots$ are omitted and
conjectured that infinitely many primes are; Booker~\cite{Booker2012}
proved the infinitude (an elementary proof is in
Pollack--Trevi\~no~\cite{PollackTrevino2014}).  Nothing comparable is known
for the least-prime-factor rule, and MC is open.

This paper does not make progress on MC.  It records, with machine-checked
proofs, what one can say about the orbit of~$2$ from the structure of~$T$,
organised in four groups.

\paragraph{Reformulations (\S\ref{sec:walk}).}  Modulo~$q$ the walk
$w_q(n)=P_n\bmod q$ satisfies $w_q(n{+}1)=w_q(n)\,m_q(n)$ with multiplier
$m_q(n)=a(n{+}1)\bmod q$, and $q\mid P_n+1$ iff $w_q(n)=-1$.  The
algebraic precondition---that the multipliers generate the group---is
free.  MC is then equivalent to each of: cofinal hitting of~$-1$; a
single hit past the sieve gap; $o(N)$ walk character sums (CCSB);
\emph{conditional multiplier equidistribution} (CME): along the orbit,
for a missing~$q$, the multiplier residues are equidistributed on each
fibre of the walk position.  All directions are machine-checked.  We
insist on the word \emph{equivalent}: every one of these statements is
quantified over primes that never appear, so MC implies it vacuously.
The chain $\CME\Rightarrow\CCSB\Rightarrow\mathrm{MC}$ therefore reduces
MC to nothing weaker; its content is the direction that turns a hitting
problem into an equidistribution problem, the shape analytic number
theory knows how to attack.

\paragraph{The composite floor (\S\ref{sec:floor}).}  The limit
$C=\lim_N \log P_N/2^{N}$ exists (a doubly-exponential-recurrence fact of
Aho--Sloane type~\cite{AhoSloane1973}, here made exact by a defect
telescope), and $C=0$ iff infinitely many Euclid numbers $P_n+1$ are
composite~$\Cinf$.  On the complementary branch---all $P_n+1$ prime from
some stage on---the recursion closes into $P_{n+1}=P_n(P_n+1)$, Sylvester's
sequence started at $P_N$, and the walk mod~$q$ becomes the autonomous
map $w\mapsto w^2+w$, which never reaches~$-1$ when $q\equiv2\pmod3$; so
MC implies $\Cinf$, i.e.\ MC implies that a doubly exponential
Sylvester-type sequence is not all prime---an open statement of the same
shape as ``infinitely many Fermat numbers are composite''.  Weak forms of
MC (reciprocal divergence RD, and a smallness statement (S)) bottom out at
$\Cinf$ as well; a second, unscaled invariant $\rho$ refines the ladder.

\paragraph{Population laws (\S\ref{sec:population}).}  The density of
$\{m:\lpf m=p\}$ is $w_p=p^{-1}\prod_{r<p}(1-1/r)=c(p)-c(p{+}1)$, so among
$q$-rough integers the class $\lpf\equiv a\pmod q$ has density
$\sum_{p\equiv a}w_p$: a convergent series dominated by its first terms.
Fixed-modulus equidistribution of the least prime factor---an assumption
that a previous version of this project placed beneath CME---is therefore
equivalent to a family of identities that Dirichlet's theorem does not
touch, and it is false.  What is true is the \emph{bag-conditioned} law:
on the progression $m\equiv1\pmod P$ that contains every Euclid number
with accumulator~$P$, the least prime outside the bag has relative density
exactly $1/q$.  Unconditionally, for each prime~$q$ almost every squarefree
seed reaches~$q$ somewhere in its factor tree; the analytic input,
$\sum_{p\equiv a}1/p=\infty$, is itself machine-checked from Mathlib's
$L(1,\chi)\ne0$, and Karamata's Tauberian theorem gives Mertens' theorem in
progressions.

\paragraph{Min versus max (\S\ref{sec:minmax}).}  Cox--van der Poorten's
omission of~$5$ for the largest-factor variant is machine-checked (a
mod-$12$ argument).  For the least-factor rule, capture $\lpf N=q$ is a
residue-class condition on~$N$, so by Dirichlet and CRT no propagating
congruence invariant---fixed or growing modulus---can block a missing
prime.  This separates the two sequences at the level of congruence
obstructions; it says nothing about their truth values.

\paragraph{The formalization (\S\ref{sec:lean}).}  Everything above is
proved in Lean~4 over Mathlib (\EMFileCount~files, $\EMLineCount$~lines,
no \code{sorry}, no \code{native\_decide}, standard axioms only).  Each
theorem carries a link {\color{leanink}\sffamily\footnotesize$\checkmark$}
to the checked declaration; the statements printed here are the ones
checked.  A companion technical report (the repository's long paper)
covers the variants, the spectral routes, a catalogue of \DEentries~dead
ends, and the methodology.

\section{The residue walk and the equivalences}
\label{sec:walk}

\subsection{The walk}

Fix a prime~$q$.  Set $w_q(n)=P_n\bmod q$ and $m_q(n)=a(n{+}1)\bmod q$, so
that $w_q(n{+}1)=w_q(n)\,m_q(n)$.

\begin{proposition}[Walk--divisibility bridge;
\lean{EM/Group/Core.lean}{walkZ\_eq\_neg\_one\_iff}]
$w_q(n)=-1$ iff $q\mid P_n+1$.
\end{proposition}

Consequently $q$ appears at stage $n{+}1$ iff $q\mid P_n+1$ and no prime
below~$q$ divides $P_n+1$; once $q$ has appeared, $w_q\equiv0$ forever.  If
$q$ never appears, $w_q(n)$ and $m_q(n)$ are units for all~$n$
(\lean{EM/Core/Conjectures.lean}{prime\_not\_in\_seq\_not\_dvd\_prod}).

\begin{theorem}[Generation is free;
\lean{EM/Equidist/Bootstrap.lean}{prime\_residue\_escape}]
For every prime~$q$, the residues of the primes below~$q$ generate
$(\ZZ/q\ZZ)^\times$.  Hence, along an inductive proof of MC in which all
primes below~$q$ have already appeared, the multipliers $m_q(n)$ generate
the group.
\end{theorem}

\begin{theorem}[Single hit;
\lean{EM/Equidist/Bootstrap.lean}{single\_hit\_implies\_mc}]
Suppose that for every prime~$q$ that never appears, if all primes below
$q$ appear and the multipliers generate $(\ZZ/q\ZZ)^\times$, then
$q\mid P_n+1$ for some~$n$ past the stage where all primes below~$q$ have
appeared.  Then MC holds.
\end{theorem}

The proof is strong induction on~$q$: past the sieve gap no prime below
$q$ can divide $P_n+1$, so a single hit gives $a(n{+}1)=q$.

\subsection{Character sums and conditional equidistribution}

For a nontrivial character $\chi$ of $(\ZZ/q\ZZ)^\times$ and a missing~$q$,
write $S_\chi(N)=\sum_{n<N}\chi(w_q(n))$.

\begin{definition}
\begin{itemize}[nosep]
\item \defname{ComplexCharSumBound} (CCSB): for every missing prime~$q$ and
every nontrivial~$\chi$, $S_\chi(N)=o(N)$
(\lean{EM/Equidist/FourierB.lean}{ComplexCharSumBound}).
\item \defname{ConditionalMultiplierEquidist} (CME): for every missing
prime~$q$, every nontrivial~$\chi$ and every position $c\in(\ZZ/q\ZZ)^\times$,
\[
  \sum_{\substack{n<N\\ w_q(n)=c}}\chi\bigl(m_q(n)\bigr)=o(N)
\]
(\lean{EM/CME/Reduction.lean}{ConditionalMultiplierEquidist}).
\end{itemize}
\end{definition}

The fibre sums are normalised by~$N$, not by the fibre count, so
zero-density fibres satisfy CME automatically and on positive-density
fibres CME says the multiplier is $\chi$-equidistributed.

\begin{theorem}[\lean{EM/CME/Reduction.lean}{cme\_implies\_ccsb},
\lean{EM/Equidist/SelfCorrecting.lean}{complex\_csb\_mc'},
\lean{EM/CME/Reduction.lean}{cme\_implies\_mc}]
$\CME\Rightarrow\CCSB\Rightarrow\mathrm{MC}$.
\end{theorem}

\begin{proof}[Sketch]
Telescoping, $\chi(w(n{+}1))=\chi(w(n))\chi(m(n))$, gives
$S_\chi(N)=\sum_{n<N}\chi(w(n))\chi(m(n))+O(1)=\sum_c\chi(c)\sum_{w(n)=c}\chi(m(n))+O(1)$,
so CME gives CCSB.  Fourier inversion on $(\ZZ/q\ZZ)^\times$ turns CCSB into
a lower bound $\#\{n<N:w_q(n)=-1\}\ge N/(q-1)-o(N)$
(\lean{EM/Equidist/FourierB.lean}{walk\_hit\_count\_fourier\_step}), i.e.\
cofinal hitting, and the single-hit theorem finishes.
\end{proof}

\begin{theorem}[The walk layer is a set of reformulations;
\lean{EM/CME/Equivalences.lean}{walk\_layer\_equivalences}]
Each of the following is equivalent to MC: cofinal hitting for every
missing prime (HH); the same conditioned on generation (DH); the
single-hit hypothesis (SHH); walk equidistribution; CME; CCSB; the
one-horizon window gain.  Threshold forms (multi-modular CCSB,
subquadratic visit energy, visit equidistribution, self-correcting drift,
decorrelation) are implied by MC.
\end{theorem}

\begin{proof}
Each statement is of the form ``for every prime~$q$ that never appears,
$\Phi(q)$''.  Under MC no prime never appears, so each holds vacuously;
the converses are the theorems above.
\end{proof}

We state this because earlier drafts of this project called CME ``the
sharpest sufficient condition known'' and treated CME, CCSB, DH as
strictly ordered; as global propositions they are one proposition.  The
hierarchy is one of proof obligations \emph{at a fixed missing~$q$}: CME
asks less per prime than CCSB (a conditional statement about the
multiplier rather than a bound on the walk sum), and CME $\Rightarrow$ MC
even has a two-line proof---on a fibre~$c$ visited $\ge N/(q-1)$ times the
class $-c^{-1}$ has zero mass if $q$ never appears, so some $\chi$ has fibre
sum of size $\ge N/(q-1)^2$.  A hypothesis with content independent of MC
would have to speak about the residues of $a(n)$ \emph{without}
conditioning on missingness, uniformly in~$q$; the population layer that
tried to supply one is discussed, and refuted, in \S\ref{sec:headdom}.

\subsection{Two illustrations in $(\ZZ/5\ZZ)^\times$}

Two periodic multiplier sequences $A=(2,2,3,3)^\infty$ and
$B=(2,3,2,3)^\infty$ in $(\ZZ/5\ZZ)^\times$ have identical multiplier
frequencies in every window of four steps; the walk of $A$ from~$1$ is
$1,2,4,2,1,\dots$ and hits $-1=4$, the walk of $B$ is $1,2,1,2,\dots$ and
never does
(\lean{EM/Meta/OrbitBarrier.lean}{population\_does\_not\_determine\_hitting}).
With $\chi(2)=i$, the multiplier sums of $B$ are bounded while
$|S_\chi(N)|=\Theta(N)$
(\lean{EM/Meta/OrbitBarrier.lean}{mult\_cancel\_not\_walk\_cancel}).  These
are elementary facts about ordering---first-order multiplier statistics do
not determine a walk, and cancellation of multiplier sums does not force
cancellation of walk sums---and they are exactly why every statement in
this section is conditional on the walk position.  We do not claim more
for them.

\section{The composite floor}
\label{sec:floor}

\subsection{The growth constant and the defect telescope}

Let $L(n)=\log P_n$ and $\delta(n)=L(n)-\log a(n{+}1)$, the amount by which
the selected prime falls short of the maximal choice $P_n$; then
$L(n{+}1)=2L(n)-\delta(n)$
(\lean{EM/Population/DefectTelescope.lean}{logProd\_succ\_eq\_two\_mul\_sub\_defect}).
If $P_n+1$ is prime the multiplier is $P_n+1$ and $\delta(n)\le0$; if it is
composite, $\lpf(P_n+1)^2\le P_n+1$ and $\delta(n)\ge\tfrac12L(n)-o(1)$
(\lean{EM/Population/DefectTelescope.lean}{defect\_gap}).

\begin{theorem}[\lean{EM/Population/DefectTelescope.lean}{tendsto\_normLog},
\lean{EM/Population/DefectTelescope.lean}{infinitelyManyComposite\_iff\_growthConstant\_eq\_zero},
\lean{EM/Population/DefectTelescope.lean}{growthConstant\_pos\_iff\_perpetualPrimality}]
$C:=\lim_{N}L(N)/2^{N}$ exists and is $\ge0$.  Moreover $C=0$ iff
infinitely many Euclid numbers $P_n+1$ are composite $\Cinf$, and $C>0$ iff
$P_n+1$ is prime for all large~$n$.  Quantitatively,
$L(N)\le(\log2+\tfrac13)(\tfrac34)^{\#\{n<N:\,P_n+1\text{ composite}\}}2^{N}$
(\lean{EM/Population/DefectTelescope.lean}{logProd\_le\_pow\_mul}).
\end{theorem}

\begin{proof}[Sketch]
Iterating the telescope, $L(N)/2^{N}=L(0)-\sum_{n<N}2^{-(n+1)}\delta(n)$.
The negative part of $\delta$ is summable after weighting: at a prime stage
$-\delta(n)=\log(1+1/P_n)\le1/P_n\le4^{-n}$, so $L(N)/2^{N}+\tfrac13 4^{-N}$
is antitone and bounded below by~$0$, hence converges; this is the
Aho--Sloane argument~\cite{AhoSloane1973} with the error made explicit.
At a composite stage $\delta(n)\ge\tfrac12L(n)-o(1)$, so
$L(n{+}1)\le\tfrac32L(n)+o(1)$, i.e.\ each composite stage multiplies the
normalised quantity by $\tfrac34$ up to a summable error, giving the
displayed bound and $C=0$ under $\Cinf$; conversely if only finitely many
stages are composite then $\delta(n)\le0$ eventually and $L(N)/2^{N}$ is
eventually nondecreasing and positive.
\end{proof}  Nothing here uses the
seed: for every $m\ge2$ the same $C(m)$ exists along the orbit of~$m$, and
the tail identity makes it a semiconjugacy to doubling,
$C(Tm)=2C(m)$ (\lean{EM/Population/SeededGrowth.lean}{sgrowth\_T}); the
statement ``$\Cinf$ from every seed'' is exactly $\{C>0\}=\emptyset$
(\lean{EM/Population/SeededGrowth.lean}{mixedDiversity\_iff\_sgrowth\_zero}),
and for every fixed threshold~$N$ the seeds whose orbit is a tower of
primes from stage~$N$ on form a set of density zero
(\lean{EM/Population/GrowthDensity.lean}{hasDensityZero\_perpetual}; primes
have density zero by the elementary sieve, no prime number theorem is
used).  We do \emph{not} claim that $\{C>0\}$ itself is null: natural
density is not countably subadditive.

\subsection{Sylvester towers, and MC $\Rightarrow\Cinf$}

On the branch $C>0$ we have $a(n{+}1)=P_n+1$ for $n\ge N$, hence
$P_{n+1}=P_n(P_n+1)$: the sequence $P_N,P_{N+1},\dots$ is Sylvester's
recursion $s_{k+1}=s_k^2+s_k$ started at $P_N$, and $\Cinf$ says exactly
that no such tower of Euclid numbers is all prime
(\lean{EM/Population/SylvesterTower.lean}{infinitelyManyComposite\_iff\_tower\_composite}).
On that branch the walk mod~$q$ is autonomous, $w\mapsto w^2+w$
(\lean{EM/Population/AutonomousBranch.lean}{perpetual\_walkZ\_eq\_ffAutonomousMap}),
and reaches $-1$ only through a root of $w^2+w+1$, which exists in $\FF_q$
only for $q\equiv1\pmod3$ (and $q=3$).

\begin{theorem}[\lean{EM/Population/AutonomousBranch.lean}{eventually\_prime\_implies\_not\_mullin},
\lean{EM/Population/CompositeFloor.lean}{infinitelyManyComposite\_of\_mullin}]
If $P_n+1$ is prime for all large~$n$ then every prime $q\equiv2\pmod3$ is
missing.  Hence $\mathrm{MC}\Rightarrow\Cinf$.
\end{theorem}

Thus MC implies that a doubly exponential Sylvester-type sequence
contains infinitely many composites---a statement of the same shape as
``infinitely many Fermat numbers $2^{2^n}+1$ are composite'', which is
open; whether Sylvester's own sequence contains infinitely many primes is
open as well~\cite{Vardi1991}, though the primes dividing it have density
zero~\cite{Odoni1985}.  We draw the modest conclusion: any proof of MC
must, along the way, rule out a Sylvester-type prime tower.  Every
conjecture is at least as hard as its open consequences, so this is a
statement of shape, not a lower bound on difficulty in any technical
sense.

\subsection{Weak forms and the invariant $\rho$}

Let RD be $\sum_k1/a(k)=\infty$ and (S) the statement that
$\lpf(P_n+1)<2^{n-c}$ for infinitely many~$n$, for every~$c$.  Then
$\mathrm{MC}\Rightarrow\RD\Rightarrow(\mathrm S)\Rightarrow\Cinf$, all
implications proved
(\lean{EM/Population/WeakMullin.lean}{mc\_implies\_rd},
\lean{EM/Population/CompositeFloor.lean}{exists\_small\_minFac\_of\_reciprocalDivergence}).
The invariant that is \emph{not} scaled by $T$ is
\[
  \rho(m)=\liminf_{n}\frac{\log\lpf(T^{n}m+1)}{\log T^{n}m},
  \qquad \rho(Tm)=\rho(m)
\]
(\lean{EM/Population/RelativeSize.lean}{rho\_T}).

\begin{theorem}[\lean{EM/Population/RelativeSize.lean}{rho\_dichotomy},
\lean{EM/Population/RelativeSize.lean}{rho\_two\_eq\_zero\_of\_rd},
\lean{EM/Population/RelativeSize.lean}{rho\_two\_le\_half\_iff}]
For every seed, $\rho\in[0,\tfrac12]\cup\{1\}$, with $\rho=1$ iff $C>0$
and $\rho\le\tfrac12$ iff $C=0$.  For the orbit of~$2$:
$\mathrm{MC}\Rightarrow\RD\Rightarrow\rho(2)=0\Rightarrow\rho(2)\le\tfrac12
\Leftrightarrow\Cinf$.
\end{theorem}

\begin{proof}[Sketch]
At a prime stage the multiplier is $P+1$, so the ratio is $1+O(1/P)$; at a
composite stage $\lpf(P+1)^2\le P+1$, so the ratio is $\le\tfrac12+O(1/P)$;
this gives the dichotomy and both equivalences.  If $\rho>\varepsilon>0$
then eventually $\lpf(T^{n}m+1)\ge(T^{n}m)^{\varepsilon}\ge2^{\varepsilon n}$ and
$\sum1/\lpf(T^{n}m+1)$ converges, so RD forces $\rho=0$.
\end{proof}

Whether $\rho(2)=0$ is strictly between RD and $\Cinf$ is open; formal
profiles show neither implication reverses for arbitrary prime sequences,
and separating them for genuine orbits is an anatomy question of the same
kind as $\Cinf$.  Along the orbit the growth coordinate and the walk
coordinate are coupled through one bit only: at a prime stage the residue
is forced ($m_q(n)=w_q(n)+1$), and at a composite stage every walk position
$w\ne-1$, residue~$a$ and size bound~$K$ are realised together by some
squarefree seed
(\lean{EM/Population/SizeResidueDecoupling.lean}{exists\_seed\_composite\_residue\_size}).
This is why the growth projection carries the floor and nothing above it.

\section{Population laws}
\label{sec:population}

\subsection{Head domination}
\label{sec:headdom}

For a prime~$p$ let $N_p=\prod_{r<p}r$ and $c(n)=\prod_{r<n,\ r\text{ prime}}(1-1/r)$.

\begin{theorem}[\lean{EM/Population/HeadDomination.lean}{card\_minFac\_eq\_ge},
\lean{EM/Population/HeadDomination.lean}{card\_minFac\_eq\_le},
\lean{EM/Population/HeadDomination.lean}{w\_eq\_cfun\_sub},
\lean{EM/Population/HeadDomination.lean}{hasSum\_wq}]
$\lfloor X/(pN_p)\rfloor\varphi(N_p)\le\#\{m\le X:\lpf m=p\}\le
(\lfloor X/(pN_p)\rfloor+1)\varphi(N_p)$, so $\{m:\lpf m=p\}$ has density
$w_p=c(p)/p=c(p)-c(p{+}1)$, and $\sum_{p>q}w_p=c(q{+}1)$.
\end{theorem}

\begin{proof}[Sketch]
$\lpf m=p$ iff $p\mid m$ and $\gcd(m,N_p)=1$; the first condition is one
class mod~$p$, the second is $\varphi(N_p)$ classes mod~$N_p$, and
$\gcd(p,N_p)=1$, so the set is a union of $\varphi(N_p)$ classes mod
$pN_p$; counting classes in blocks gives the two bounds.  The telescoping
$c(p)/p=c(p)-c(p{+}1)$ is $c(p{+}1)=c(p)(1-1/p)$, and $c(n)\to0$ because
$\sum_p1/p=\infty$; hence $\sum_{p>q}w_p=c(q{+}1)$.
\end{proof}

\begin{corollary}[\lean{EM/Population/HeadDomination.lean}{tendsto\_classCount\_div},
\lean{EM/Population/HeadDomination.lean}{roughLPFEquidist\_iff}]
Among $q$-rough integers ($\lpf m>q$) the class $\lpf m\equiv a\pmod q$ has
density $\sum_{p\equiv a\,(q)}w_p$.  Hence ``$\lpf\bmod q$ is equidistributed
among the coprime classes of $q$-rough integers'' is equivalent to
$\sum_{p\equiv a}w_p=c(q{+}1)/(q-1)$ for every unit class~$a$.
\end{corollary}

These identities are false: the series is dominated by its first terms,
the class of the least prime $p_0>q$ receives more than its share (about
twice, for large~$q$), and Dirichlet's theorem, which controls the
massless tail, cannot restore the balance
(\lean{EM/Population/HeadDomination.lean}{not\_roughLPFEquidist\_of\_head}
is the finite-head criterion).  We state the falsity as a fact about the
identities and, by a standing rule of this project, do not include a
numerical certificate for a specific~$q$; the shifted version (Euclid
numbers) and its uniform-conductor version fail identically, and the
statement over all integers is refuted outright
(\lean{EM/Equidist/SieveTransfer.lean}{not\_genericLPFEquidist}: $\lpf n\equiv1
\pmod3$ forces $\gcd(n,6)=1$).  A previous version of this project placed
such an equidistribution statement beneath CME, attributing it to Alladi;
Alladi's theorem~\cite{Alladi1977} is the M\"obius-weighted duality
$-\sum_{n\le x,\,\lpf n\equiv a}\mu(n)\sim x/\varphi(q)$, which is a
different statement.

\subsection{The bag-conditioned law}

Every Euclid number $P_n+1$ lies in the progression $1\pmod{P_n}$, and its
least prime factor is a prime outside the bag.

\begin{theorem}[\lean{EM/Population/BagConditionedLaw.lean}{tendsto\_bagClass\_div\_ap},
\lean{EM/Population/BagConditionedLaw.lean}{bagWeight\_least\_missing}]
For $P\ge1$ and a prime $p\nmid P$,
\[
  \lim_{X\to\infty}\frac{\#\{m\le X:\ m\equiv1\ (P),\ \lpf m=p\}}
  {\#\{m\le X:\ m\equiv1\ (P)\}}
  =\frac1p\prod_{r<p,\ r\nmid P}\Bigl(1-\frac1r\Bigr).
\]
In particular, if $q$ is the least prime not dividing~$P$, the relative
density of $\{\lpf m=q\}$ on the progression is exactly $1/q$.
\end{theorem}

\begin{proof}[Sketch]
For $m\equiv1\pmod P$ every prime dividing $P$ is prime to~$m$, so
$\lpf m=p$ iff $p\mid m$ and $\gcd(m,N')=1$ with
$N'=\prod_{r<p,\ r\nmid P}r$.  The conditions $m\equiv1\ (P)$ and
$m\equiv0\ (p)$ combine by CRT into one class $c_0$ mod $pP$, and along the
progression $m=c_0+pP\,t$ the coprimality to $N'$ holds for exactly
$\varphi(N')$ values of $t$ in each block of $N'$ consecutive integers,
because $t\mapsto(pP\,t+c_0)\bmod N'$ is a bijection of $\ZZ/N'$
(\lean{EM/ForMathlib/CoprimeAffineBlock.lean}{card\_coprime\_affine\_block}).
Divide by the count $\sim X/P$ of the progression.
\end{proof}

\noindent This is
Shanks' ``probability $1/q$'' heuristic~\cite{Shanks1991} made precise at
the population level; what the orbit inherits from it remains heuristic,
since $P_n+1$ is one member of the progression, not a random one.

\subsection{The seed-average law: almost all seeds coprime to $q$ capture $q$}
\label{sec:seed-average}

The bag-conditioned law describes one step.  Iterating it gives an
unconditional density theorem for the greedy map itself.  Run the greedy
map from an arbitrary seed: $P_m(0)=m$, $P_m(k{+}1)=P_m(k)\,p_m(k)$ with
$p_m(k)=\lpf(P_m(k)+1)$.

\begin{theorem}[\lean{EM/Population/AlmostAllDensity.lean}{almost\_all\_genmc\_density},
\lean{EM/Population/AlmostAllDensity.lean}{almost\_all\_genmc\_limsup}]
\label{thm:almost-all-capture}
For every prime~$q$ and every $\varepsilon>0$ there are a horizon~$n$ and a
threshold~$X_0$ such that for all $X\ge X_0$,
$\#\{m\le X: q\nmid m,\ p_m(j)\ne q \text{ for all } j<n\}\le\varepsilon X$.
\end{theorem}

\noindent\emph{Scope.}  This is a statement about the \emph{seed ensemble}
and says nothing about the Euclid--Mullin orbit of~$2$, which is a single
trajectory.  The prime~$q$ is \emph{fixed}: $n$ and $X_0$ depend on $q$ and
$\varepsilon$, and the simultaneous-in-$q$ form is open \emph{for natural
density}, since upper natural density is only finitely \emph{sub}additive:
finite subadditivity does give one
horizon uniform over any \emph{finite} set of primes
(\lean{EM/Population/AlmostAllDensity.lean}{finite\_simultaneous\_density}),
and a diagonal argument lets the range grow with the seed: for some
ineffective nondecreasing $Q\to\infty$, almost every seed~$m$ selects every
prime $q\le Q(m)$ coprime to it
(\lean{EM/Population/GrowingRange.lean}{seed\_range\_never\_density}).
That is as far as it goes.  It is proved on a different ensemble,
which carries a countably additive measure (next paragraph); the
two are not comparable, because the integers are null there.  The horizon
is finite; there is no
limit in~$n$, no equidistribution hypothesis, and every input is
unconditional.

\noindent\emph{The profinite packaging.}  Simultaneity in~$q$ was never a
rate problem --- summable per-$q$ failure fractions are available, by
applying the theorem at $\varepsilon_q=q^{-2}$ --- but an additivity
problem, and the repair is a countably additive ambient measure.  Take
$\Omega=\prod_{r\ \text{prime}}\ZZ/r\ZZ$ with $\mu$ the product of the
uniform measures.  Write $p_x(j)$ for the $j$-th multiplier of the greedy
orbit of a point $x\in\Omega$, defined next.  Coordinates $\ZZ/r\ZZ$, not $\ZZ_r$: the programme only
ever conditions on $m\bmod M$ with $M$ squarefree, so one residue per
prime records everything it can see, and including \emph{every} prime as a
coordinate --- $q$ among them --- is what makes the $q$-coordinate the free
CRT coordinate that the selection law exploits by excluding~$q$ from its
own modulus $M_Y=\prod_{r\le Y,\,r\ne q}r$ (movement (ii) below).
The measure enters once, as an equality: an $M$-periodic event has measure
exactly its period fraction $\#T/M$
(\lean{EM/Population/ProfiniteEnsemble.lean}{measure\_residue\_classes}).
A profinite point has no $\lpf$, so the multiplier is re-expressed
coordinatewise as the least prime whose coordinate of the current Euclid
element vanishes
(\lean{EM/Population/ProfiniteDynamics.lean}{leastVanishing}); along the
diagonal $\iota(m)=(m\bmod r)_r$ this \emph{is} the integer dynamics, with
no hypothesis beyond $m\ge1$
(\lean{EM/Population/ProfiniteDynamics.lean}{profSeq\_iota}) --- the same
process, not a lookalike.  Lifting a profinite point by CRT to an integer
agreeing with it at every prime coordinate $\le Y$ (with $q\le Y$
load-bearing, since $M_Y$ omits~$q$) puts every point that misses~$q$ into
the bad-type covering of movement (iii), so
$\mu\{x: x_q\ne0,\ p_x(j)\ne q\ \forall j\}=0$
(\lean{EM/Population/ProfiniteHeadline.lean}{measure\_missing\_eq\_zero}),
and countable additivity does the union over~$q$:
\[
  \mu\bigl\{x\in\Omega:\ \exists q\ \text{prime},\ x_q\ne0
     \ \text{and}\ p_x(j)\ne q\ \text{for all } j\bigr\}=0
\]
(\lean{EM/Population/ProfiniteHeadline.lean}{measure\_some\_prime\_missed\_eq\_zero}).
No $q$-uniform rate is used anywhere: at each fixed~$q$ the horizon goes to
infinity separately.  \emph{Scope.}  $\NN\subset\Omega$ is $\mu$-null
(\lean{EM/Population/ProfiniteEnsemble.lean}{measure\_range\_iota\_eq\_zero},
proved), so ``$\mu$-a.e.\ seed'' is not ``almost all integer seeds'' --- a
$\mu$-null set can have upper density~$1$ in~$\NN$; structurally, $\mu$-a.e.\ point
has infinitely many vanishing coordinates while an integer has finitely many
(\lean{EM/Population/ProfiniteAttractor.lean}{measure\_divisorFinite\_eq\_zero}),
and MC itself reads ``$\Prod{}(n)\to0$ in the profinite topology''
(\lean{EM/Population/ProfiniteAttractor.lean}{mc\_iff\_tendsto\_zero}) --- and transfer must go
through the finite residue-class statements plus a genuine
equidistribution input, never through ``a.e.'' alone.  Mathematically new
content: none; the finite counting chain is the mathematics and this is
packaging, which is a feature, not an apology.  Nothing follows about the
orbit of~$2$.

The proof has three movements, and forms no character sum anywhere.

\emph{(i) Capture becomes a residue condition.}  A truncated orbit is a CRT
function of the seed: if $M$ is divisible by every prime $\le y$ and the
first $n$ multipliers of $m$ stay $\le y$, then every $m'\equiv m\ (M)$ has
the same first $n$ multipliers
(\lean{EM/Population/SeedTypes.lean}{genSeq\_prefix\_of\_modEq}).  This
fails exactly at~$q$, so one runs the \emph{$q$-free} dynamics, which always
selects the least prime $\ne q$ dividing the current Euclid number; that run
is a legitimate, seed-residue-blind reference orbit.  Let $V_q(m,n)$ be the
residues mod~$q$ of its cofactors at the steps where it declined an
available~$q$.  Then for $q\nmid m'$, $m'\equiv m\ (M)$,
\[
  \exists\,j<n:\ p_{m'}(j)=q
  \iff
  m'\bmod q\ \in\ -V_q(m,n)^{-1},
  \qquad \#V_q(m,n)\le n
\]
(\lean{EM/Population/SeedCapture.lean}{captured\_iff\_mem\_visited}).
Dynamical capture is membership of one residue in an explicitly enumerated
set of size $\le n$.

\emph{(ii) An exact selection law and a deterministic compensator.}  With
$M_Y=\prod_{r\le Y,\,r\ne q}r$ --- $q$ excluded, so that the coordinate
$m\bmod q$ stays CRT-free --- the seeds of one period $[1,M_Y]$ split into
\emph{type cells}, the fibres of the first-$n$-multipliers map, each cut out
by local residue conditions
(\lean{EM/Population/SelectionLaw.lean}{mem\_cell\_iff\_local}).  Inside a
cell the proportion of seeds whose next Euclid number has no prime factor
$\le y$ other than $q$ is \emph{exactly} the roughness survival product
$S=\prod_{r\le y,\,r\ne q}(1-\rho_r)$, $\rho_r=1/\#B_r$ the reciprocal of a
box of surviving unit residues
(\lean{EM/Population/SelectionLaw.lean}{selection\_law}).  This is Shanks'
heuristic made exact \emph{and} conditional on the whole type so far; the
$1/q$ law above is its $n=0$ case.  The survival products are not summable,
by a conservation law: each active prime shrinks its own box by one every
time it declines, so it pays at most $H_{r-1}$ in total
(\lean{EM/Population/LargeStepRoughness.lean}{charge\_sum\_le\_harmonic}), and
$\sum_{r<2n}H_{r-1}=O(n)$.  Hence, on \emph{every} nondegenerate $Y$-bounded
type path with $\log Y\le n^2$,
$\sum_{k<n}S_k\ge (c_1/2)n$ with $c_1=e^{-250}$ absolute
(\lean{EM/Population/LargeStepRoughness.lean}{pathwise\_compensator}).  This
bound is deterministic and per-path: no probability, no ergodicity.  Feeding
it, with the selection law as the conditional counting hypothesis, into an
abstract finite-tree exponential supermartingale
(\lean{EM/Population/TreeChernoff.lean}{chernoff\_quarter\_local}; Mathlib
only, no measure theory) and into Lemma~D --- large steps land in any
prescribed class coprime to~$q$ with proportion $\ge\kappa=e^{-128}/(16\varphi(q))$,
from Karamata at window scale
(\lean{EM/Population/LemmaDBox.lean}{lemma\_D\_z}) --- gives Theorem~C: the
good uncaptured seeds are at most $M_Y\exp(-\tfrac38\kappa((c_1/2)n-K_0))$
(\lean{EM/Population/TheoremC.lean}{theorem\_C}), whence an $\varepsilon$
bound on one period
(\lean{EM/Population/AlmostAllGenMC.lean}{almost\_all\_genmc}).

\emph{(iii) Period fraction to natural density.}  This step is not a
formality.  A period fraction is a diagonal count and natural density is a
product count, and no inequality relates them: an $M$-periodic set can have
diagonal count~$0$ over one period and product count $(q{-}1)M/q$.  Transfer
therefore requires that \emph{every} clause be a function of $m\bmod M_Y$,
and two are not --- $q\nmid m$ and ``uncaptured'', which by (i) is a
condition on the fibre coordinate $m\bmod q$, absent from $M_Y$.  The repair
costs no constant: the Chernoff argument uses capture only through the
readoff ``visited set full $\Rightarrow$ every unit residue captured'' and
through a capture criterion already stated for a general seed of the fibre,
so both clauses weaken to ``\emph{some} seed of the residue fibre of~$m$ is
coprime to~$q$ and uncaptured''
(\lean{EM/Population/FiberTheoremC.lean}{FiberGood}), which is
$M_Y$-periodic, and Theorem~C survives verbatim
(\lean{EM/Population/FiberTheoremC.lean}{theorem\_C\_fiber}).  All three bad
types are then fibre-measurable and jointly rare
(\lean{EM/Population/TypeBadSmall.lean}{type\_bad\_small}), and a periodic
covering of density $\varepsilon/2$ gives upper density $\varepsilon$ on
$[1,X]$
(\lean{EM/ForMathlib/PeriodicDensity.lean}{eventually\_density\_le}), the
factor~$2$ absorbing the incomplete final block.  The truncation~$Y$ appears
in neither the hypotheses nor the conclusion of
Theorem~\ref{thm:almost-all-capture}.

\subsection{Almost-all hitting in the factor tree}

Replace the deterministic choice by the \emph{factor tree}: from a
squarefree seed $m$, branch to $m\cdot p$ for every prime $p\mid m+1$.  Say
$m$ \emph{reaches}~$q$ if some node of the tree is $\equiv-1\pmod q$, i.e.\
$q$ divides some Euclid number along some path.

\begin{theorem}[\lean{EM/Ensemble/UnconditionalPSCD.lean}{almost\_all\_mixed\_hitting\_unconditional},
\lean{EM/Ensemble/UnconditionalPSCD.lean}{almost\_all\_mixed\_hitting\_diagonal}]
For each prime~$q$, the squarefree seeds coprime to~$q$ that do not
reach~$q$ have density zero among squarefree integers; and there is
$B(X)\to\infty$ such that almost every squarefree seed $\le X$ reaches
every prime $\le B(X)$.
\end{theorem}

\begin{proof}[Sketch]
One step suffices: $m\cdot p\equiv-1\pmod q$ for a prime $p\mid m+1$ iff
$p\equiv-m^{-1}\pmod q$.  The seeds $m\le X$ for which $m+1$ has no prime
factor $\le z$ in a given unit class~$a$ number at most
$X\prod_{p\le z,\ p\equiv a}(1-1/p)+O(2^{\pi(z)})$ by an elementary sieve,
and $\prod_{p\le z,\ p\equiv a}(1-1/p)\to0$ because
$\sum_{p\equiv a}1/p=\infty$
(\lean{EM/IK/DirichletDensity.lean}{prime\_reciprocal\_class\_divergent},
from Mathlib's $L(1,\chi)\ne0$); a weak unconditional lower bound
$\#\{m\le X\text{ squarefree}\}\ge X/4$ absorbs the constants.  The
diagonal statement chooses $B(X)$ growing slowly enough that the finitely
many bad densities are still small.
\end{proof}  The
theorem is modest---one step already suffices, since $m+1$ has a prime
factor in the required class for almost all~$m$---and it is about the
tree, not the greedy path.

\subsection{Karamata, and Mertens in progressions}

\begin{theorem}[\lean{EM/IK/Karamata.lean}{karamata}]
Let $c_n\ge0$, $\sum_n c_n n^{-s}<\infty$ for every $s>0$, and
$s\sum_n c_n n^{-s}\to C$ as $s\to0^+$.  Then $\sum_{n\le x}c_n\sim C\log x$.
\end{theorem}

\begin{proof}[Sketch]
Write $F(s)=\sum_nc_nn^{-s}$ and $y=n^{-s}$.  The hypothesis gives
$sF(s)\to C$, and for every $k\ge0$, $s\sum_nc_nn^{-s(k+1)}\to C/(k+1)
=C\int_0^1y^{k}\,dy$; by linearity the same holds with $y^{k}$ replaced by
any polynomial.  Sandwich the function $y\mapsto y^{-1}\mathbf 1_{[e^{-1},1]}(y)$
between continuous ramps and, by Weierstrass, between polynomials whose
integrals differ by less than~$\eta$; since all $c_n\ge0$ the sums are
monotone in the test function, and $\sum_nc_nn^{-s}\cdot n^{s}\mathbf 1_{n\le e^{1/s}}
=\sum_{n\le e^{1/s}}c_n$ is squeezed to $C/s\,(1+O(\eta))$.  Put $x=e^{1/s}$.
\end{proof}

Applied to $c_n=\Lambda(n)\mathbf 1_{n\equiv a\,(q)}$, whose Dirichlet series
has the pole supplied by Mathlib's $L(1,\chi)\ne0$, this gives Mertens'
theorem in progressions,
$\sum_{p\le x,\ p\equiv a}\log p/p\sim\log x/\varphi(q)$ and
$\sum_{p\le x,\ p\equiv a}1/p\sim\log\log x/\varphi(q)$
(\lean{EM/IK/Karamata.lean}{weightedPNTinAP\_asymp\_proved},
\lean{EM/IK/Karamata.lean}{primesEquidistInAP\_asymp\_proved}).  Mathlib has
neither a Tauberian theorem nor Mertens' theorem, so these are offered as
infrastructure.\footnote{Mertens' first theorem is not currently available
in Mathlib (as of \texttt{v4.33.0}); a pull request upstreaming it is open
at the time of writing~\cite{Mathlib41394}.  The explicit two-sided bounds
used in \S\ref{sec:seed-average}
(\lean{EM/Population/MertensLower.lean}{mertens\_lower}, and a windowed
$\log\log$ lower bound) are therefore proved here from scratch by an
elementary route, with deliberately crude constants.  No priority is
claimed: sharper formalizations exist elsewhere --- Eberl and Paulson's
Isabelle/HOL AFP entry gives
$\mathfrak{M}(n)-\ln n\in(-1-9/\pi^2,\ \ln4]$~\cite{EberlPaulsonPNT}, and the
\texttt{PrimeNumberTheoremAnd} Lean project proves
$|\sum_{p\le x}\log p/p-\log x|\le\log4+4$~\cite{PNTplus}.}

\section{Min versus max}
\label{sec:minmax}

Let $b(0)=2$, $b(n{+}1)=\gpf(b(0)\cdots b(n)+1)$ be the largest-prime-factor
variant.

\begin{theorem}[Cox--van der Poorten, machine-checked;
\lean{EM/Obstruction/MaxVariant.lean}{five\_missing}]
$5$ never appears in $(b(n))$.
\end{theorem}

\begin{proof}[Sketch]
From stage~$1$ on the accumulator is divisible by $2$ and $3$ and is
$\equiv2\pmod4$, so the candidate $N=P_n+1$ is odd, prime to~$3$, and
$\equiv3\pmod4$.  If $\gpf N=5$ then $N$ has no prime factor other than
$5$, so $N=5^c\equiv1\pmod4$, a contradiction
(\lean{EM/Obstruction/MaxVariant.lean}{maxFac\_ne\_five\_of\_mod\_four}).
The invariant ``$P\equiv6\pmod{12}$'' propagates and blocks~$5$.
\end{proof}

The reason this works for $\gpf$ and not for $\lpf$ is a one-line
asymmetry.

\begin{proposition}
$\lpf N=q$ iff $q\mid N$ and $r\nmid N$ for every prime $r<q$: a
residue-class condition on $N$ modulo $q\prod_{r<q}r$.  The condition
$\gpf N=q$ is not a residue-class condition modulo any fixed integer.
\end{proposition}

Consequently, for the least-factor rule, an ``invariant'' that sees the
candidate $N=P_n+1$ only through its residue class modulo some $m$ can
never certify that $q$ is missed: Dirichlet supplies candidates with
$\lpf N=q$ in every admissible class.  Formally, call $S\subseteq\ZZ/m$
\emph{propagating} if it is closed under every transition
$r\mapsto r\cdot\lpf(N)$ with $N\equiv r+1\pmod m$, and say $S$
\emph{blocks}~$q$ if no state of $S$ forces capture of~$q$.

\begin{theorem}[\lean{EM/Obstruction/NoInvariant.lean}{no\_cvdp\_obstruction},
\lean{EM/Reciprocity/NoReciprocityInvariant.lean}{no\_reciprocity\_induction\_proof}]
For a missing prime~$q$ and any modulus $m\ne0$, no propagating set
containing the tail of the orbit blocks~$q$; the same holds when the
modulus is allowed to grow with the orbit (which covers Jacobi-symbol
invariants against $P_n$).
\end{theorem}

\begin{proof}[Sketch]
Since $q$ is missing, from some stage the accumulator is prime to~$m$
and its residue $r$ satisfies: $r+1$ is a unit mod~$m$ (or, at worst, the
class $r+1$ contains integers prime to $m$).  By CRT and Dirichlet there is
$N\equiv r+1\pmod m$ with $\lpf N=q$ (a residue-class condition on~$N$
modulo $q\prod_{r'<q}r'$, compatible with the class mod~$m$ after removing
common factors, since $q$ is prime to the accumulator).  So the transition
$r\mapsto r\cdot q$ is available from every tail state, and a propagating
set containing the tail contains a state that forces capture; it cannot
block~$q$.  The growing-modulus version runs the same argument at
$\Pi_n=8mP_n$, where a Jacobi symbol against $P_n$ is a function of the
residue.
\end{proof}

The transition relation is deliberately over-approximated ($N$ ranges over
the whole residue class), which makes the non-existence statement
\emph{weaker} than one for the true dynamics; the honest reading is about
proofs: an omission argument for $(a(n))$ whose only knowledge of the
candidate is congruential cannot succeed, so any such proof must use the
anatomy of $P_n+1$---its large prime factors, exactly what $\lpf$
discards.  Nothing here constrains proofs outside this fragment, and
Booker's argument for the largest-factor sequence is a finite-window
reductio using smoothness, not an invariant of this shape.

\section{The formalization}
\label{sec:lean}

All results are formalized in Lean~4 against Mathlib \texttt{v4.33.0}
(\EMFileCount~files, $\EMLineCount$~lines).  The verified set is the import
closure of the root module: no \code{sorry}, no user axioms, no
\code{native\_decide}; every published declaration depends only on
\code{propext}, \code{Classical.choice}, \code{Quot.sound} (checked by a
script over the registry).  The definitions of the sequence
(\code{Mullin.seq}, \code{Mullin.prod}) use a home-grown constructive
\code{minFac} bridged to Mathlib's, and \code{MullinConjecture} is
literally ``every prime is some \code{seq n}''.  Each theorem environment
above names the checked declaration; the statements printed here were
written from the Lean statements, and where they abbreviate (e.g.\ the
$o(N)$ notation for an explicit $\forall\varepsilon\,\exists N_0$) the Lean
text is the reference.  Published results and open hypotheses are tagged
\code{@[publish]}/\code{@[open\_point]} and exported to a machine-readable
registry (\code{registry/}); the open points are MC and its equivalents,
the threshold hypotheses, and the floor ladder.

The repository's long document is a technical report: the variants
(factor tree, two-point, $\varepsilon$-perturbed, mixed, function field,
with placeholders marked), the spectral and variance routes, the
obstruction calculus, a complete catalogue of \DEentries~documented dead
ends with the reason each fails, and the human--AI methodology.  Nothing
in this paper depends on it.

\bigskip
\noindent\textbf{Acknowledgements.}  The Lean code was written with AI
assistance (Claude, Anthropic); the mathematics, the choice of what to
prove, and the errors are the author's.

\begin{thebibliography}{9}
\bibitem{Mullin1963} A.\,A.~Mullin, Recursive function theory (a modern look at a Euclidean idea), \emph{Bull.\ Amer.\ Math.\ Soc.}\ 69 (1963), 737.
\bibitem{Shanks1991} D.~Shanks, Euclid's primes, \emph{Bull.\ Inst.\ Combin.\ Appl.}\ 1 (1991), 33--36.
\bibitem{CoxVdP1968} C.\,D.~Cox and A.\,J.~van der Poorten, On a sequence of prime numbers, \emph{J.\ Austral.\ Math.\ Soc.}\ 8 (1968), 571--574.
\bibitem{Booker2012} A.\,R.~Booker, On Mullin's second sequence of primes, \emph{Integers} 12A (2012), A4.
\bibitem{PollackTrevino2014} P.~Pollack and E.~Trevi\~no, The primes that Euclid forgot, \emph{Amer.\ Math.\ Monthly} 121 (2014), 433--437.
\bibitem{AhoSloane1973} A.\,V.~Aho and N.\,J.\,A.~Sloane, Some doubly exponential sequences, \emph{Fibonacci Quart.}\ 11 (1973), 429--437.
\bibitem{Vardi1991} I.~Vardi, Are all Euclid numbers squarefree?, in \emph{Computational Recreations in Mathematica}, Addison-Wesley, 1991, 82--89.
\bibitem{Odoni1985} R.\,W.\,K.~Odoni, On the prime divisors of the sequence $w_{n+1}=1+w_1\cdots w_n$, \emph{J.\ London Math.\ Soc.\ (2)} 32 (1985), 1--11.
\bibitem{Alladi1977} K.~Alladi, Duality between prime factors and an application to the prime number theorem for arithmetic progressions, \emph{J.\ Number Theory} 9 (1977), 436--451.
\bibitem{EberlPaulsonPNT} M.~Eberl and L.\,C.~Paulson, The prime number theorem, \emph{Archive of Formal Proofs}, 2018; Isabelle/HOL theory \texttt{Mertens\_Theorems}. \url{https://www.isa-afp.org/entries/Prime_Number_Theorem.html}
\bibitem{PNTplus} A.~Kontorovich, T.~Tao, et al., \emph{PrimeNumberTheoremAnd}, Lean~4 project, 2024--2026; file \texttt{PrimeNumberTheoremAnd/IEANTN/Mertens.lean}. \url{https://github.com/AlexKontorovich/PrimeNumberTheoremAnd}
\bibitem{Mathlib41394} T.~Tao, \texttt{feat(NumberTheory/Mertens)}: the Mertens theorems, mathlib4 pull request \#41394, July 2026 (open at the time of writing). \url{https://github.com/leanprover-community/mathlib4/pull/41394}
\end{thebibliography}

\end{document}
